\documentclass[12pt, letterpaper]{article}

%% Build from the paper/ directory:
%%   cd paper && pdflatex main.tex
%% or use the root build system:
%%   make paper          (from repo root)
%%   ./build.sh paper    (POSIX)
%%   build.cmd paper     (Windows cmd)

% packages.tex
\usepackage{amsmath, amssymb, amsthm}
\usepackage{graphicx}
\usepackage{xcolor}
\usepackage[
    hidelinks,
    pdftitle={Tidal Ionization of Atomic Hydrogen and the Quantum Roche Limit},
    pdfauthor={Jack Menendez},
    pdfsubject={Quantum Roche limit; gradient ionization; Paper III of the A=1 Discrete Causal Lattice series},
    pdfkeywords={A=1 Discrete Causal Lattice, quantum Roche limit, tidal ionization, gradient ionization, Arnold tongue, clock density, atomic ISCO, 21cm HI, accretion disk}
]{hyperref}
\usepackage{booktabs}
\usepackage{array}
\usepackage{geometry}
\usepackage{adjustbox}
\usepackage{longtable}
\usepackage{setspace}
\usepackage{caption}
\usepackage{subcaption}
\usepackage{float}
\onehalfspacing

\input{macros/commands}

%% Page geometry (also set here as a fallback in case the package path fails)
\geometry{top=1in, bottom=1in, left=1in, right=1in}

%% ---------------------------------------------------------------
%% TITLE PAGE
%% ---------------------------------------------------------------

\title{Tidal Ionization of Atomic Hydrogen and\\
       the Quantum Roche Limit%
       \thanks{Version 0.1-DRAFT.  Paper~III of the A=1 Discrete
       Causal Lattice series; takes up follow-on item~\#9
       (\emph{Tidal Ionization Mass}) of Paper~I's
       \texttt{notes/follow\_on\_implications.md} catalogue.
       Builds on \emph{Geometry First}
       (Paper~I,
       \href{https://doi.org/10.5281/zenodo.20078529}{DOI:
       10.5281/zenodo.20078529}) and is orthogonal to Paper~II
       (\emph{Geometry Forces Physics},
       \href{https://doi.org/10.5281/zenodo.20240736}{DOI:
       10.5281/zenodo.20240736}).
       Pre-release draft; DOI for this paper to be assigned at
       v1.0 deposit.}}
\author{Jack Menendez\\
        \small Eugene, OR\\
        \small ORCID: \href{https://orcid.org/0009-0003-1166-307X}{0009-0003-1166-307X}\\
        \small \href{mailto:jack@geometryinducedphysics.org}{jack@geometryinducedphysics.org}\\
        \small Web: \href{https://geometryinducedphysics.org/}{geometryinducedphysics.org}}
\date{\today}

\begin{document}

\maketitle

\begin{abstract}
    % abstract.tex

The hydrogen atom in the A=1 Discrete Causal Lattice framework
(Paper~I~\cite{menendez2026geometry}) is a joint Arnold-tongue
attractor between the proton and electron sessions, lock-in
occurring when their orbital frequency ratio matches the lattice's
$\omega \cdot R_1 = \pi/3$ resonance to within the tongue width
$\Delta\omega_\text{tongue}$.  An external mass $M$ at distance $d$
imposes a differential clock-density gradient across the atom of
order $\Delta\omega_\text{tidal} \sim M \cdot R_1 / d^3$.  When the
tidal detuning exceeds the tongue width, the joint resonance breaks
and the orbit dissolves -- gradient ionization, a third channel
beyond thermal and pressure ionization.  This paper derives the
critical condition $M > \Delta\omega_\text{tongue} \cdot d^3 / R_1$
analytically -- the \emph{quantum Roche limit}, sharing the $d^3$
scaling of the classical Roche limit with $\Delta\omega_\text{tongue}$
playing the role of the classical body density -- and confirms it
numerically via the inherited \texttt{exp\_18\_tidal\_ionization.py}
g-scan, which measures the orbital escape time $T_\text{escape}(g)$
as the imposed gradient $g = M / d^2$ is swept.  The framework's
parameter-free $M_\text{min}(d)$ ionization curve gives a
temperature-independent inner edge for neutral hydrogen near compact
objects, testable against 21\,cm HI absorption morphology and
accretion-disk Fe K$\alpha$ line profiles.  The paper closes with the
atomic ISCO -- the innermost stable atomic orbit, an analogue of the
geodesic ISCO at $r = 6GM/c^2$, below which no neutral matter can
exist regardless of temperature.

\end{abstract}

%% ---------------------------------------------------------------
%% FRONT-MATTER NOTICES (after abstract per author preference)
%% ---------------------------------------------------------------

\medskip
\noindent\textbf{Keywords.}  A=1 Discrete Causal Lattice; quantum
Roche limit; tidal ionization; gradient ionization; Arnold tongue;
clock density; atomic ISCO; 21\,cm HI; compact-object atmospheres.

\medskip
\noindent\textbf{License.}  Paper text and figures: Creative Commons
Attribution 4.0 (CC BY 4.0).  Source code in the companion
repository: MIT License.

\medskip
\noindent\textbf{Data availability.}  Numerical experiments, data
files, and the \LaTeX{} source for this paper are publicly available
at \nolinkurl{github.com/JackDMenendez/dcl-paper-03-tidal-ionization}.
The load-bearing numerical experiment
\nolinkurl{exp_18_tidal_ionization.py} was introduced in Paper~I
(\texttt{src/experiments/} of the upstream dcl
repo)~\cite{menendez2026geometry} and is reproduced here verbatim
with a Provenance block; Paper~III's contribution is the
analytical $M_\text{min}(d)$ quantum Roche limit derivation, the
$d^3$ scaling identification, the numerical confirmation produced by
running the inherited experiment, and the observational predictions
for 21\,cm HI inner-edge morphology near compact objects.

%% ---------------------------------------------------------------
%% RELEASE NOTES (only at release time)
%% ---------------------------------------------------------------
%% \bigskip
%% \noindent\textbf{Release notes for vX.Y (since vX.Y-1).}
%% \begin{itemize}
%% \item Item 1.
%% \end{itemize}

%% \emph{Status legend used throughout.}  \texttt{PASS} / \texttt{PART}
%% / \texttt{STUB} / \texttt{FAIL} -- see Table~\ref{tab:audit}.

\newpage
\tableofcontents
\newpage

%% ============================================================
%% MAIN
%% ============================================================
\section{Introduction}
% introduction.tex

\label{sec:introduction}

\placeholder{Phase~2 content: write the orientation paragraphs that
frame the quantum Roche limit as a third ionization channel
distinct from thermal and pressure ionization.  Headline arc:
(i)~standard astrophysics recognises two ionization channels
(thermal Saha, pressure-degenerate); (ii)~the A=1 framework
predicts a third (gradient detuning of the joint Arnold-tongue
lock-in); (iii)~the gradient-ionization condition reduces to the
quantum Roche limit $M > \Delta\omega_\text{tongue} \cdot d^3 / R_1$,
with the $d^3$ scaling identifying it as the quantum analogue of the
classical tidal-disruption radius~\cite{roche1849}; (iv)~the
prediction is parameter-free and testable against 21\,cm HI inner-
edge morphology near compact objects.}

\subsection{Scope}
\label{subsec:scope}

In scope for this paper:

\begin{itemize}
\item Analytical derivation of $M_\text{min}(d)$ from the
      Arnold-tongue-width / gradient-detuning condition.
\item Identification of the $d^3$ scaling as the quantum analogue
      of the classical Roche limit.
\item Numerical confirmation via \texttt{exp\_18\_tidal\_ionization.py}
      (g-scan: $T_\text{escape}(g)$ decreases monotonically with $g$
      and drops sharply near $g_\text{crit} = \Delta\omega_\text{tongue}
      / R_1$).
\item The atomic ISCO: the innermost stable atomic orbit below
      which no neutral matter can exist independent of temperature.
\item Observational predictions for 21\,cm HI inner-edge morphology
      near compact objects and accretion-disk Fe K$\alpha$ line
      profiles.
\end{itemize}

Out of scope (deferred to companion papers):

\begin{itemize}
\item The \emph{full} cosmological-recombination application
      (follow-on \#2 of Paper~I's catalogue).
\item The Arnold-tongue-compression-at-higher-$n$ analysis
      ($\Delta\omega_\text{tongue}(n) \sim 1/n^2$,
      $R_n \sim n^2$, threshold $\sim 1/n^4$) is referenced from
      Paper~I's \texttt{notes/plasma\_as\_gravitational\_ionization.md}
      but not re-derived here -- the present paper focuses on the
      $n=1$ ground state.
\item The relativistic-infall regime (follow-on \#10 of Paper~I's
      catalogue, Case~3) where time dilation partially cancels
      gradient ionization for radially infalling atoms.
\end{itemize}

Table~\ref{tab:audit} records the audit status of every claim in
the paper.  At v0.1-DRAFT, the analytical and numerical-confirmation
rows are \texttt{STUB} (paper-text derivation pending) and
\texttt{PART} (script in place, runtime evidence pending the first
run); the audit table will be updated as Phase~2 lands.

%% Audit table embedded here so the reader meets it before any
%% claim that depends on a status entry.
% ============================================================
% audit_table.tex -- Paper III: Tidal Ionization / Quantum Roche Limit
%
% A "system audit" table mapping the paper's claims to the evidence
% backing each one.  This is *the* canonical source of truth for which
% claims are PASS / PART / STUB / FAIL --- prose elsewhere in the paper
% should not contradict it.
%
% Status semantics:
%   PASS  -- derivation complete or experiment confirms the claim to
%            stated precision.
%   PART  -- mechanism demonstrated, full quantitative match pending;
%            specific gaps documented in the evidence column.
%   STUB  -- analytical result with no numerical confirmation yet,
%            or planned experiment / paper-text derivation not yet
%            written.
%   FAIL  -- tested and disconfirmed.
% ============================================================

\label{tab:audit}

{\scriptsize
\setlength{\tabcolsep}{4pt}
\renewcommand{\arraystretch}{0.95}
\begin{longtable}{@{}p{2.8cm}p{3.6cm}p{3.6cm}p{4.0cm}p{1.0cm}@{}}
\caption{\small Audit for Paper~III.  Inherited rows are evidenced by
Paper~I material (analytical derivations in
\texttt{notes/plasma\_as\_gravitational\_ionization.md}; numerical
machinery in \texttt{exp\_18\_tidal\_ionization.py}, reproduced
verbatim in this repo).  Paper~III-specific rows -- the explicit
$M_\text{min}(d)$ derivation, the $d^3$ scaling identification, the
g-scan run results, and the observational predictions -- start
\texttt{STUB}/\texttt{PART} in v0.1 and flip as Phase~2 lands.}
\label{tab:audit_caption} \\
\hline
\textbf{Claim} & \textbf{Mechanism} & \textbf{Continuum / formal limit} & \textbf{Evidence} & \textbf{Status} \\
\hline
\endfirsthead

\multicolumn{5}{l}{\textit{(Table~\ref{tab:audit}, continued from previous page)}} \\
\hline
\textbf{Claim} & \textbf{Mechanism} & \textbf{Continuum / formal limit} & \textbf{Evidence} & \textbf{Status} \\
\hline
\endhead

\hline
\multicolumn{5}{r}{\textit{(continued on next page)}} \\
\endfoot

\hline
\endlastfoot

% ── Inherited from Paper I ──────────────────────────────────────────
\multicolumn{5}{l}{\textit{Inherited from Paper~I~\cite{menendez2026geometry}.}} \\
\hline

Hydrogen as joint two-session Arnold-tongue lock-in
    & Bipartite tick rule + clock-density mass; lock-in at
      $\omega \cdot R_1 = \pi/3$
    & Standard hydrogen spectrum in the continuum limit
    & \texttt{exp\_12} (Paper~I), $k_\text{min} = 0.0970$ vs
      $k_\text{Bohr} = 0.0971$ to 4 sig figs
    & \texttt{PASS} \\

Gradient detuning breaks the joint resonance
    & $\omega_\text{eff}(x) = \omega \cdot \bar{\rho} /
      \rho_\text{clock}(x)$; differential detuning across $R_1$
      exceeds $\Delta\omega_\text{tongue}$
    & Three ionization channels (thermal Saha + pressure-degenerate
      + gradient detuning); the third is new in this framework
    & \texttt{notes/plasma\_as\_gravitational\_ionization.md}
      (analytical)
    & \texttt{PASS} \\

% ── Paper III-specific ──────────────────────────────────────────────
\hline
\multicolumn{5}{l}{\textit{Paper~III-specific claims (this paper).}} \\
\hline

Quantum Roche limit $M_\text{min}(d) \cdot R_1 / d^3 =
\Delta\omega_\text{tongue}$
    & Substitute $g = M / d^2$ into the gradient-detuning condition;
      solve for $M$ at threshold
    & $M_\text{min}(d) = \Delta\omega_\text{tongue} \cdot d^3 / R_1$;
      same $d^3$ scaling as classical Roche limit~\cite{roche1849}
      with $\Delta\omega_\text{tongue}$ replacing body density
    & Paper-text derivation pending (\S\,planned)
    & \texttt{STUB} \\

$d^3$ scaling is the quantum-mechanical analogue of the classical
Roche limit
    & Both derive from $|\nabla F| \cdot \ell / F = $ const at
      threshold, where $\ell$ is the bound-state size and $F$ is
      the strength of the binding force; the framework derives the
      quantum case from the Arnold-tongue width
    & The two limits are the same identity seen from two levels of
      description; in the lattice they are the same thing
    & Paper-text observation pending (\S\,planned)
    & \texttt{STUB} \\

Three ionization channels (thermal, pressure, gradient)
    & Saha (kinetic) + Pauli (degenerate) + Arnold-tongue
      detuning (gradient); the third is new in the framework
    & Standard astrophysics recognises the first two; the third
      becomes dominant near compact objects
    & \texttt{notes/plasma\_as\_gravitational\_ionization.md};
      paper-text consolidation pending
    & \texttt{PART} \\

Numerical confirmation:
$T_\text{escape}(g)$ decreases monotonically with $g$, sharp drop
near $g_\text{crit}$
    & g-scan in \texttt{exp\_18}: apply linear gradient
      $V_\text{tidal} = g \cdot (x - x_\text{centre})$ from
      $t = 0$; measure tick at which $r_\text{pdf} > 2 R_1$ for
      $N$ consecutive windows
    & $g_\text{crit} \approx \Delta\omega_\text{tongue} / R_1$;
      Paper~I's free-particle tongue width gives
      $g_\text{crit} \approx 0.016$ per node
    & \texttt{exp\_18\_tidal\_ionization.py}: script in place
      (Paper~I copy, reproduced here); runtime evidence pending
      first run from this repo
    & \texttt{STUB} \\

Tidal displacement diagnostic:
electron and proton move in opposite directions under gradient
(not isotropically as for thermal disruption)
    & x-projection of CoM relative to system CoM should show
      $x_e \cdot x_p < 0$ for gradient ionization, isotropic
      scatter for energy-injection disruption
    & Distinguishes gradient ionization from
      kinetic-energy-injection mechanisms
    & \texttt{exp\_18} diagnostic columns
      (\texttt{x\_e\_proj}, \texttt{x\_p\_proj}); runtime evidence
      pending
    & \texttt{STUB} \\

Atomic ISCO: innermost stable atomic orbit
$r_\text{atomic\_ISCO}(M)$
    & Largest $d$ at which gradient ionization always succeeds
      for the $n=1$ ground state, independent of temperature;
      analogue of geodesic ISCO at $r = 6 GM / c^2$
    & Innermost neutral-hydrogen radius around compact objects;
      structural prediction independent of thermal modelling
    & Paper-text derivation pending (\S\,planned)
    & \texttt{STUB} \\

Observational signature:
21\,cm HI inner-edge morphology near compact objects
    & Threshold $M_\text{min}(d)$ gives a temperature-independent
      inner boundary for neutral hydrogen
    & 21\,cm HI absorption / emission inner edge near neutron stars
      and active galactic nuclei
    & Paper-text comparison to existing 21\,cm HI surveys pending
      (\S\,planned)
    & \texttt{STUB} \\

Observational signature:
Fe K$\alpha$ line-profile inner edge in accretion disks
    & Atomic ISCO gives an inner edge for Fe atoms supporting
      K$\alpha$ emission; below the atomic ISCO, Fe is fully
      ionized regardless of temperature
    & Inner edge of Fe K$\alpha$ emission near accretion disks
      (XMM-Newton / NICER profiles)
    & Paper-text comparison pending (\S\,planned)
    & \texttt{STUB} \\

\hline
\end{longtable}
}



%% Planned sections for v1.0 (Phase 2 content, drafted as we go):
%%   \section{The Arnold-tongue lock-in mechanism}
%%   \section{Gradient detuning and the critical condition}
%%   \section{The quantum Roche limit and its $d^3$ scaling}
%%   \section{Numerical confirmation (exp\_18)}
%%   \section{Observational predictions}

\section{Conclusion}
% conclusion.tex

\placeholder{Conclusion to be written in Phase~2.  Headline arc:
(i)~the quantum Roche limit is a parameter-free prediction with the
same $d^3$ scaling as the classical Roche limit, derived from a
single condition on the Arnold-tongue width; (ii)~the atomic ISCO
is the discrete-substrate analogue of the geodesic ISCO, predicting
an innermost neutral-hydrogen radius around compact objects;
(iii)~the prediction is testable now against existing 21\,cm HI and
Fe K$\alpha$ data without new instrumentation; (iv)~the relativistic-
infall regime (follow-on~\#10 of Paper~I) is the natural Paper~IV
extension where time dilation partially cancels gradient ionization
for radially infalling gas.}


\section*{Acknowledgements}
% acknowledgements.tex

\placeholder{Acknowledge collaborators, reviewers, and tools.  The
inherited \texttt{exp\_18\_tidal\_ionization.py} script and the
\texttt{notes/plasma\_as\_gravitational\_ionization.md} theoretical
analysis were introduced in Paper~I of the A=1 Discrete Causal
Lattice series.  Acknowledgement of arXiv endorsers (if any) and
Zenodo community curators (the \texttt{a1-discrete-causal-lattice},
\texttt{fbt-framework}, and \texttt{osqgr} communities at minimum)
belongs here once the paper is deposited.}


%% ============================================================
%% APPENDICES
%% ============================================================
\appendix

\section{Code and Data}
% code_and_data.tex
\label{sec:code_and_data}

\subsection{Repository layout}

\begin{itemize}
\item \texttt{src/experiments/exp\_18\_tidal\_ionization.py}
  -- the load-bearing g-scan experiment.  Reproduced verbatim from
  Paper~I's
  \texttt{src/experiments/exp\_18\_tidal\_ionization.py}~\cite{menendez2026geometry};
  see the Provenance block in the script's docstring.  Sweeps
  gradient $g$ from $0$ (control) through $\sim\!5\times
  g_\text{crit}$, measures $T_\text{escape}(g)$ via the windowed
  $r_\text{pdf}$ diagnostic, and records the tidal-displacement
  signature ($x_e, x_p$ projections relative to system CoM).
\item \texttt{src/experiments/exp\_18\_tidal\_ionization.md} --
  Paper~III's expansion of Paper~I's stub: parameter table, expected
  runtime, PASS/PART/FAIL criteria for each audit-table row this
  experiment evidences.
\item \texttt{notes/plasma\_as\_gravitational\_ionization.md} --
  inherited from Paper~I with a Provenance prefix.  The
  Arnold-tongue / clock-density-gradient theoretical apparatus this
  paper builds on.
\item \texttt{notes/quantum\_roche\_limit.md} -- new in this paper.
  Synthesises Paper~I's gradient-ionization argument with the $M/d^3$
  Roche-scaling observation from Paper~I's follow-on catalogue
  entry~\#9, and outlines the $M_\text{min}(d)$ derivation that lands
  in the paper's main body.
\item \texttt{data/exp\_18\_tidal\_ionization.npy} -- g-scan output:
  rows of $(g, \text{escaped}, T_\text{escape}, x_e\_proj,
  x_p\_proj, r_\text{pdf,first\_window})$.  Written by the
  experiment script; tracked in git once the run produces it.
\item \texttt{data/exp\_18\_rpdf\_trajectories.npy} -- full
  $r_\text{pdf}(t)$ trajectories per $g$ value, for the trajectory
  figures.
\end{itemize}

\subsection{Engine}

The experiment script imports framework primitives
(\texttt{OctahedralLattice}, \texttt{CausalSession},
\texttt{enforce\_unity\_spinor}) from \texttt{dcl\_core.core} --
the continuous-amplitude submodule of the shared \texttt{dcl\_core}
package
(\href{https://github.com/JackDMenendez/dcl-core}{\nolinkurl{github.com/JackDMenendez/dcl-core}}),
which contains a verbatim port of Paper~I's
original \texttt{src/core/}.  The dependency is declared in
\texttt{virtual-env-requirements.txt} and installed automatically
by \texttt{setup.cmd} / \texttt{setup.sh}; no separate engine
provisioning is required.  At Paper~III's v1.0 release, the
\texttt{dcl\_core} dependency will be pinned to a specific
\texttt{@vX.Y.Z} tag (and that tag's Zenodo deposit DOI) so that
the released paper binds to an immutable engine version.


\section{Reproducibility Quickstart}
% reproducibility.tex
\label{sec:reproducibility}

\begin{itemize}
\item \textbf{Environment.}  Python~3 with \texttt{numpy},
      \texttt{scipy}, \texttt{matplotlib}, \texttt{pytest}.
      Exact dependency list in
      \texttt{virtual-env-requirements.txt}.  The \texttt{setup.sh}
      / \texttt{setup.cmd} scripts create a \texttt{.venv} and
      install everything.

\item \textbf{Engine.}  The experiment depends on the A=1 framework
      engine via the shared \texttt{dcl\_core} package
      (\texttt{dcl\_core.core.OctahedralLattice},
      \texttt{dcl\_core.core.CausalSession}, \dots).
      \texttt{dcl\_core} is declared in
      \texttt{virtual-env-requirements.txt} and installed by
      \texttt{setup.cmd} / \texttt{setup.sh} automatically; no
      separate engine setup step is required.  See \texttt{CLAUDE.md}'s
      \emph{Engine} section for the architectural rationale and the
      pinning strategy at release time.

\item \textbf{Build the paper.}  From a fresh clone, with the venv
      activated:
      \begin{verbatim}
./build.sh paper        # POSIX / MSYS2 UCRT64 on Windows
build.cmd paper         # Windows cmd / PowerShell
      \end{verbatim}
      runs \texttt{pdflatex} three times and \texttt{bibtex} once,
      producing \texttt{paper/main.pdf}.

\item \textbf{Run the experiment.}  Once the engine is provisioned:
      \begin{verbatim}
python -u src/experiments/exp_18_tidal_ionization.py \
    > data/exp_18_tidal_ionization.log 2>&1
      \end{verbatim}
      Wall-clock runtime: a few hours on a modern laptop (7
      gradient values $\times$ up to 2500 ticks each, with each
      tick a 3D wavefunction step on a $65^3$ grid).  Expected
      output: monotonic $T_\text{escape}(g)$ decrease and an
      opposite-sign tidal-displacement signature
      ($x_e\_proj \cdot x_p\_proj < 0$) for $g > 0$.

\item \textbf{Master audit roll-up.}  From the repo root:
      \begin{verbatim}
python audit_universe.py
      \end{verbatim}
      reports the declared status of every row in
      \texttt{paper/sections/audit\_table.tex} and parses the
      latest \texttt{data/exp\_18*.log} for the actual cached
      \texttt{PASS}/\texttt{FAIL} marker.

\item \textbf{Where to look next.}  Start with
      \texttt{notes/quantum\_roche\_limit.md} for the analytical
      derivation; then
      \texttt{notes/plasma\_as\_gravitational\_ionization.md}
      (inherited from Paper~I) for the broader gradient-ionization
      theory; then the experiment's \texttt{.md} companion for
      runtime parameters and audit criteria.
\end{itemize}


%% ============================================================
\bibliographystyle{unsrt}
\bibliography{paper-bib/references}

\end{document}
